%! Author = chtompki
%! Date = 1/14/26
\documentclass[12pt]{amsart}
% Default margins are too wide all the way around. I reset them here
\setlength{\hoffset}{-.75in}
\setlength{\textwidth}{6.5in}
\setlength{\voffset}{-.5in}
\setlength{\textheight}{9.0in}
\usepackage{amsmath}
\usepackage{leftindex}
\usepackage{amssymb}
\usepackage{tikz}
\usepackage{amsthm}
\usepackage{enumitem}
\usepackage{setspace}

\theoremstyle{definition}
\newtheorem{definition}{Definition}

\newenvironment{solution}
{\begin{proof}[Solution]}
{\end{proof}}

\title{Combinatorics\\Notes January 15, 2026}
%\author{Rob Tompkins}
\date{\today}

\begin{document}
    \maketitle
    \date{\today}
    \begin{onehalfspace}
        \underline{Counting}
        \begin{quote}
            ``Those who do not count; do not count!''
        \end{quote}
        Our combinatorial proofs will mainly rely upon the following principles
        \begin{enumerate}
            \item The Addition rule,
            \item The Multiplication rule,
            \item The Permutation rule (with or without repeating),
            \item The Combination rule (with or without repeating),
            \item The Distribution rule (twelve fold method),
            \item Inclusion and Exclusion,
            \item Generating Functions (ordinary and exponential),
            \item Recurring Relations (linear and non-linear),
            \item The Pigeonhole Principle (both simple and generalized) \emph{important},
            \item Other miscellaneous methods.
        \end{enumerate}
        Suppose we have boxes for distributing items
        \begin{figure}[!ht]
            \centering
            \begin{tikzpicture}
                \draw (-4,1) -- (-4, 0) -- (-2, 0) -- (-2, 1);
                \draw (-1, 1) -- (-1, 0) -- (1, 0) -- (1, 1);
                \draw (2, 1) -- (2, 0) -- (4, 0) -- (4, 1);
            \end{tikzpicture}
            \caption{Three boxes}
        \end{figure}
        We'd want to quickly ask some questions about our boxes:
        \begin{itemize}
            \item are they labled?
            \item are the items indistinguishable or have types?
            \item Define a bijection for capturing the mechanics of the problem.
        \end{itemize}
        \emph{Note.} Maybe use the ``Sterling'' or ``Burr'' Numbers.
        To solve this problem we'd probably use rules 1, 2, 3, or 4 from above.

        Before we get into the principles at hand let us note one definition.
        \begin{definition}
            A \emph{partition} of a group of $r$ identical objects divides the group into a collection
            of (unordered) subsets of various sizes. Analogously, we define a partition of the integer $r$ to be a
            collection of positive integers who's sum is $r$. Normally we write this collection as a sum and list the
            integers in increasing order.
        \end{definition}
        For example, the seven partions of the integer 5 are:
        \begin{eqnarray*}
            1 + 1 + 1 + 1 + 1  \qquad  1 + 1 + 1 + 2 \qquad 1 + 1 + 3 \\
            1 + 2 + 2 \qquad 1 + 4 \qquad 2 + 3 \qquad 5
        \end{eqnarray*}

        \begin{definition}
            \emph{The Addition Principle.} Suppose a set $S$ is partitioned into pairwise disjoint parts
             $S_1, S_2, \ldots S_m$, the number of the parts is given by adding the numbers obtained
            \begin{displaymath}
                |S| = \sum_{i=1}^m |s_i| = |s_1| + |s_2| + \ldots + |s_m|
            \end{displaymath}
        \end{definition}

        Consider the following example where we look at a situation where the subsets or parts are not necessarily disjoint
        \begin{displaymath}
            S = S_1 \cup S_2 \cup S_3 \cup S_4,
        \end{displaymath}
        but we have
        \begin{eqnarray*}
            |S| & = & |S_1| + |S_2| + |S_3| + |S_4| - |S_1 \cap S_2| - |S_1 \cap S_3| - |S_1 \cap S_4| \\
                &  &  - |S_2 \cap S_3 | - |S_2 \cap S_4| - |S_3 \cap S_4| \\
                &  &  + | S_1 \cap S_2 \cap S_3| + |S_1 \cap S_2 \cap S_4| + |S_1 \cap S_3 \cap S_4| + |S_2 \cap S_3 \cap S_4| \\
                &  &  - |S_1 \cap S_2 \cap S_3 \cap S_4|
        \end{eqnarray*}

        \begin{definition}
            Let $S$ be a st of ordered $n$-tuples $(a_1, a_2, \ldots, a_n)$ of objects where the element $a_j$ comes
            from a set of size $n_j$ then the size of $S$ is given by
            \begin{displaymath}
                |S| = \prod_{i=1}^n n_i
            \end{displaymath}
        \end{definition}

        \noindent\textbf{Example.} Determine the nujber of integers that are the factors of $N = 3^4 \cdot 5^2 \cdot 11^7 \cdot 13^8$.
        We simply add one to the exponents and multiply them together giving us
        \begin{displaymath}
            5 \times 3 \times 8 \times 9  = 1080
        \end{displaymath}

        \begin{definition}
            \emph{Permutations.} (with no repition and order matters). \\
            $\leftindex_n P_r$ is the permutation of $n$ elements taking $r$ at a time
            \begin{displaymath}
                n \times n-1 \times n-2 \times \cdots \times n-r + 1 \times \frac{(n-r)(n-r-1)\cdots (1)}{(n-r)(n-r-1)\cdots(1)} = \frac{n!}{(n-r)!}
            \end{displaymath}
        \end{definition}

        \noindent\textbf{Example.} Consider $\{a,b,c,d\}$ let's look at $\leftindex_4 P_3 = 4 \cdot 3 \cdot 2 = 24$
        \begin{center}
            \begin{tabular}{ c|c|c|c }
                abc & bac & cab & dab \\
                abd & bca & cad & dac \\
                acd & bad & cba & dba \\
                acb & bda & cbd & dbc \\
                adb & bcd & cda & dca \\
                adc & bdc & cdb & dcb
            \end{tabular}
        \end{center}
        and if we drop the ordering then we have
        \begin{displaymath}
            \leftindex_n C_r = \frac{\leftindex_n P_r}{r!} = \frac{n!}{(n-r)!r!} = \binom{n}{r}.
        \end{displaymath}

        \noindent\textbf{Example.} How many ways are there to rearrange the word \emph{mississippi}?
        using $2!$ for the $p$'s, $4!$ for the $s$'s, $4!$ for the $i$'s and $1!$ for the $m$ we count it as
        \begin{displaymath}
            \frac{11!}{2!4!4!2!}
        \end{displaymath}
        In how many ways can we arrange the letters of \emph{mississippi} so that the vowels are together? Well how about
        we just take all six of the $i$'s and block them together calling them $I$. Then, we are merely concerned with
        the word $mIsssspp$, and that counts out as:
        \begin{displaymath}
            \frac{8!}{4!2!}
        \end{displaymath}

        \noindent\textbf{Example.} Bagel Shop. We have three types of bagels plain $P$, everything $E$, and
        cinnamon-rasin $R$. The combination for selecting 12 bagels follows as
        \begin{displaymath}
            P + E + R = 12
        \end{displaymath}
        how may non-negative integer solutions does the equation have. For example the following are solutions
        \begin{eqnarray*}
            (6, 6, 0) \\
            (4, 4, 4) \\
            (3, 3, 6)
        \end{eqnarray*}
        Also we can represent the bagles in the following pattern
        \begin{displaymath}
            \underbrace{bb}_{P} \Bigg\vert \underbrace{bbbbbb}_{E} \Bigg\vert \underbrace{bbbb}_{R}
        \end{displaymath}
        or even more generally
        \begin{displaymath}
            \vert\vert bbbbbb \ldots
        \end{displaymath}
        Regardless when we count these suckers up we get
        \begin{eqnarray*}
            \binom{12 + 3 - 1}{12} & = & \binom{12 + 3 - 1}{3 - 1} \\
            \binom{14}{12} = \binom{14}{2}
        \end{eqnarray*}

        \emph{Important.} The number of ways to select $r$ (order) from $n$ (varietes) with repitition is
        \begin{displaymath}
            \binom{n + r - 1}{r} = \binom{n + r - 1}{n-1}.
        \end{displaymath}

    \end{onehalfspace}
\end{document}
